\documentclass[10pt,a4paper]{amsart}
\usepackage{amsmath,amssymb,amsthm,graphicx}
% Minimal class-compatibility preamble for a reader excerpt of
%   Bhimarthi Ravinder, "Demazure Modules, Chari-Venkatesh Modules and Fusion Products",
%   SIGMA 10 (2014), 110, DOI 10.3842/SIGMA.2014.110.
%
% The author's pinned TeX file declares no theorem environments of its own: they are
% supplied by the SIGMA publisher class sigma.cls, which is not installed here.  The
% block below is that class's theorem-environment block (same environment names and
% styles), copied mechanically so that the retained author text compiles unchanged.
% No mathematical macro or statement is introduced by this file.
\newtheorem{theorem}{Theorem}
\newtheorem*{theorem*}{Theorem}
\newtheorem{lemma}{Lemma}
\newtheorem*{lemma*}{Lemma}
\newtheorem{corollary}{Corollary}
\newtheorem*{corollary*}{Corollary}
\newtheorem{proposition}{Proposition}
\newtheorem*{proposition*}{Proposition}
\newtheorem{conjecture}{Conjecture}
\newtheorem*{conjecture*}{Conjecture}
{\theoremstyle{definition} \newtheorem{definition}{Definition}
\newtheorem*{definition*}{Definition}
\newtheorem{example}{Example}
\newtheorem*{example*}{Example}
\newtheorem{remark}{Remark}
\newtheorem*{remark*}{Remark}
\newtheorem{note}{Note}
\newtheorem*{note*}{Note}
}
\usepackage{cleveref}
\crefname{originalreference}{original reference}{original references}
\Crefname{originalreference}{Original reference}{Original references}
\makeatletter
\newlabel{garland}{{garland}{source}{}{sourcepointer-1}{}}
\newlabel{garland@cref}{{[originalreference][0][]garland}{[1][0][]source}{}{}{}}
\makeatother
\crefname{originalreference}{original reference}{original references}
\Crefname{originalreference}{Original reference}{Original references}
\makeatletter
\newlabel{w2}{{w2}{source}{}{sourcepointer-2}{}}
\newlabel{w2@cref}{{[originalreference][0][]w2}{[1][0][]source}{}{}{}}
\makeatother
\title{Two relations in the local Weyl module attached to a multiple of the highest root}
\author{Bhimarthi Ravinder}
\date{doi:10.3842/SIGMA.2014.110; CC BY 4.0}
\begin{document}
\maketitle
\noindent\textit{Editorial note.} All delivered mathematical text is the author's own, copied from the pinned article source sources/sigma-2014-110-demazure-fusion/pinned-original-source.tex; no statement or proof step was rewritten, reordered or removed.\par\medskip
\noindent\textit{Editorial note.} Preamble provenance: the author's tex declares no theorem environments of its own (they come from the SIGMA publisher class); the staged sources/sigma-2014-110-demazure-fusion/retained-preamble.tex reproduces that class's theorem-environment block verbatim and is used as the preamble, with the wrapper class and base packages supplied by the assembler.\par\medskip
\noindent\textit{Editorial note.} Bibliography provenance: neither the counted statement nor its proof contains a citation (this lane rescans the whole span before writing the card), so the card carries no bibliography and no entry was reconstructed.\par\medskip
\noindent\textit{Editorial note.} Selection provenance: this card was opened under the file-based control contract (CONTROL.md), with the mechanical pre-card gate run on its statement and proof interval before the card was written.\par\medskip
\noindent\textit{Allowed standard prerequisites (editorial).}\begin{itemize}
\item The local Weyl module W((k+1)theta) with its highest weight vector, and the relation (w2) remembering that a sufficiently high power of a negative root vector acts as zero on it. The author's section 3, where the module and the relation are set up before the counted statement; the relation is referred to in the proof at source line 367 (a navigation-pointer stub). Fixes the module in which the two relations hold and supplies the vanishing used at the end of both parts.
\item The commutator of a raising operator with a power of a lowering operator in the current algebra, and the evaluation of the coefficient appearing in the expansion. Computed inside the counted proof at source lines 341-356. Produces the explicit multiple of the desired vector which equals a vector known to vanish.
\item The commutation relation for two negative root vectors, labelled garland in the article. Invoked inside the counted proof at source line 365 (a navigation-pointer stub). Gives the second relation after specialising the parameters.
\end{itemize}\medskip
\ifdefined\hypertarget\hypertarget{sourcepointer-1}{}\fi
\noindent\textit{Source reference.} Original reference garland: sources/sigma-2014-110-demazure-fusion/pinned-original-source.tex, source lines 315--315. Navigation pointer only: the counted text refers to the item garland of this same article, which is not part of this card. The pointer carries the item's own author label as its number; its content is not reproduced here.\par\medskip
\ifdefined\hypertarget\hypertarget{sourcepointer-2}{}\fi
\noindent\textit{Source reference.} Original reference w2: sources/sigma-2014-110-demazure-fusion/pinned-original-source.tex, source lines 307--307. Navigation pointer only: the counted text refers to the item w2 of this same article, which is not part of this card. The pointer carries the item's own author label as its number; its content is not reproduced here.\par\medskip
\begin{lemma}
\label{g}
Let $k\in\mathbb{N}$.
The following relations hold in the local Weyl module $W((k+1)\theta)$:
\begin{enumerate}\itemsep=0pt
\item[$1)$] $(x^{-}_{\theta}\otimes 1)^{2k+1}\big(x^{-}_{\theta}\otimes t^{2k+1-i}\big)w_{(k+1)\theta}=0,
\qquad
\forall\,
0\leq i \leq k$;
\item[$2)$] $(x^{-}_{\theta}\otimes t^m)(x^{-}_{\theta} \otimes t^{m+1})  w_{(k+1)\theta} \in  \big\langle \big(x^{-}_{\theta}
\otimes t^{m+2}\big) w_{(k+1)\theta} \rangle,
\qquad
\forall\,
m\geq k$.
\end{enumerate}
\end{lemma}

\begin{proof}
To prove part (1), consider $(x^{+}_{\theta}\!\otimes t^{2k{+}1{-}i}) (x^{-}_{\theta}\otimes 1)^{2k{+}3}w_{(k{+}1)\theta}$.
Since $(x^{+}_{\theta}\!\otimes t^{2k{+}1{-}i})w_{(k{+}1)\theta}$ $=0$, we get
\begin{gather*}
\big(x^{+}_{\theta}\otimes t^{2k+1-i}\big) (x^{-}_{\theta}\otimes 1)^{2k+3}w_{(k+1)\theta}
=\big[x^{+}_{\theta}\otimes t^{2k+1-i}, (x^{-}_{\theta}\otimes 1)^{2k+3}\big] w_{(k+1)\theta}
\\
\qquad{}
=\sum\limits_{j=1}^{2k+3}  (x^{-}_{\theta}\otimes 1)^{j-1}\big(\theta^{\vee}\otimes t^{2k+1-i}\big)(x^{-}_{\theta}\otimes
1)^{2k+3-j}w_{(k+1)\theta}.
\end{gather*}
Since $(\theta^{\vee}\otimes t^{2k+1-i})w_{(k+1)\theta}=0$, we may replace $(\theta^{\vee}\otimes
t^{2k+1-i})(x^{-}_{\theta}\otimes 1)^{2k+3-j}$~by
\begin{gather*}
\big[\theta^{\vee}\otimes t^{2k+1-i}, (x^{-}_{\theta}\otimes 1)^{2k+3-j}\big]
=(-2)(2k+3-j)(x^{-}_{\theta}\otimes 1)^{2k+2-j}\big(x^{-}_{\theta}\otimes t^{2k+1-i}\big).
\end{gather*}
After simplifying, we get
\begin{gather*}
\big(x^{+}_{\theta}\otimes t^{2k+1-i}\big) (x^{-}_{\theta}\otimes 1)^{2k+3}w_{(k+1)\theta}
\\
\qquad{}
=(-1)(2k+2)(2k+3)(x^{-}_{\theta}\otimes 1)^{2k+1}\big(x^{-}_{\theta}\otimes t^{2k+1-i}\big)w_{(k+1)\theta}.
\end{gather*}
Now, using $(x^{-}_{\theta}\otimes 1)^{2k+3}w_{(k+1)\theta}=0$ in $W((k+1)\theta)$, completes the proof of part~(1).
Part (2) follows easily by putting $r=2$, $s=2m+1$ and $\alpha=\theta$ in Lemma~\ref{garland}, and using the fact that
$(x^{-}_{\theta}\otimes 1)^{2m+3}w_{(k+1)\theta}=0$, $\forall\,
m\geq k$ by~\eqref{w2}.
\end{proof}

\end{document}
